% Claim proof file for row claim_3_26 -- "AcyclicNonCollidersBlockable".
% Auto-generated stub by scaffold/solving_current_row.py at row start. A
% manager agent dedicated to writing the TeX proof fills in the body below.
% The proof must:
%   - Stay close to the lecture notes' proof if one exists (always search
%     the LN first for a `\begin{proof}` block immediately following the
%     claim; copy / lightly edit when present).
%   - Otherwise use the LN paradigm and cite prior claims by their `ref`.
%   - Be self-contained enough that a mathematician can follow it without
%     re-reading the LN.
%   - Have no `sorry`, `omitted`, or "obvious" shortcuts.
%
% Compiles standalone via the `subfiles` package.

\documentclass[main]{subfiles}

\begin{document}
% ref: claim_3_26 (proof, with statement restated for readability)
% title: AcyclicNonCollidersBlockable
\def\rowref{claim\_3\_26}
\phantomsection\label{claim_3_26}
%
% Cross-references: see claim_<ref>_statement_<title>.tex in this folder
% for the corresponding statement.

% --- statement (restated) ---------------------------------------------------
\begin{Rem}[AcyclicNonCollidersBlockable]
    If a CDMG $G$ is acyclic then all non-colliders are blockable.
    So, the partial condition for $i\sigma$-separation
    ``a blockable non-collider in $C$''
    can be simplified to ``(any) non-collider in $C$''.
    
    So in the acyclic case we can simplify the notion of $i\sigma$-separation, 
    which we will refer to as \emph{$id$-separation}.
    However, in the non-acyclic setting $id$-separation (``(any) non-collider in $C$'') and $i\sigma$-separation (``a blockable non-collider in $C$'') are clearly not equivalent.

    It turned out that in the non-acyclic case $i\sigma$-separation is the more general concept (and as said above it also captures the acyclic case equivalently well), see \cite{FM17,FM18,FM20,BFPM21}.
    We will first focus on CADMGs (acyclic) for which we can restrict ourselves to the somewhat simpler $id$-separation.
    Later, we will pick up $i\sigma$-separation again when we deal with cycles.
\end{Rem}

% --- proof ------------------------------------------------------------------
\begin{proof}
The LN block contains two mathematical assertions plus a discussion paragraph: the first sentence is \refrow{claim_3_20}, already proven; the genuinely novel content is the second sentence, the simplification of the ``blockable non-collider in $C$'' clause to ``(any) non-collider in $C$''. We prove this simplification layered through the three definitions involved -- the walk-position predicate \refrow{def_3_16}, the per-walk $\sigma$-blocking predicate \refrow{def_3_17}, and $i\sigma$-separation \refrow{def_3_18} -- and close with an acknowledgment of the LN's discussion paragraph.

\smallskip
\noindent\textit{Step 0: the two predicates agree at every walk-position.}\;
By \refrow{def_3_16}, ``$v_k$ is a blockable non-collider on $\pi$'' is defined as ``$v_k$ is a non-collider on $\pi$ \emph{and} $v_k$ is not an unblockable non-collider on $\pi$''. In particular, ``blockable non-collider'' implies ``non-collider'' by taking the first conjunct, and this direction holds in every CDMG with no acyclicity hypothesis. Conversely, by \refrow{claim_3_20}, if $G$ is acyclic then for every walk $\pi$ in $G$ and every position $v_k$ on $\pi$, ``$v_k$ is a non-collider on $\pi$'' implies ``$v_k$ is a blockable non-collider on $\pi$''. Combining the two directions, on an acyclic CDMG the predicates ``$v_k$ is a blockable non-collider on $\pi$'' and ``$v_k$ is a non-collider on $\pi$'' are extensionally equivalent at every walk-position.

\smallskip
\noindent\textit{Step 1: walk-level simplification.}\;
Fix a walk $\pi$ in the acyclic CDMG $G$ and a set $C \ins J \cup V$. Conjoining the per-position equivalence of step~0 with the position-only predicate ``$v_k \in C$'' (which is independent of the blockable/non-blockable refinement) and then existentially quantifying over the position index $k$ yields
\[
  \big( \exists k,\ v_k \text{ is a blockable non-collider on } \pi \text{ and } v_k \in C \big)
  \;\iff\;
  \big( \exists k,\ v_k \text{ is a non-collider on } \pi \text{ and } v_k \in C \big).
\]
This is the LN's simplification at the walk-existential level.

\smallskip
\noindent\textit{Step 2: $\sigma$-blocked-level simplification.}\;
By \refrow{def_3_17} (clause~(2)), $\pi$ is $C$-$\sigma$-blocked iff
\[
  \big( \exists k,\ v_k \text{ is a collider on } \pi \text{ and } v_k \notin \Anc^G(C) \big)
  \quad\text{or}\quad
  \big( \exists k,\ v_k \text{ is a blockable non-collider on } \pi \text{ and } v_k \in C \big).
\]
The collider disjunct is unaffected by the simplification; substituting the walk-level equivalence of step~1 into the second (non-collider) disjunct gives the LN's \emph{simplified} form of $\sigma$-blocking,
\[
  \big( \exists k,\ v_k \text{ is a collider on } \pi \text{ and } v_k \notin \Anc^G(C) \big)
  \quad\text{or}\quad
  \big( \exists k,\ v_k \text{ is a non-collider on } \pi \text{ and } v_k \in C \big),
\]
i.e.\ the LN's second sentence read at the per-walk predicate level.

\smallskip
\noindent\textit{Step 3: $i\sigma$-separation lift.}\;
By \refrow{def_3_18} (clause~(1)), $A \isPerp_G B \given C$ unfolds to ``for every $v \in A$, every $w \in J \cup B$, and every walk $\pi$ from $v$ to $w$ in $G$, $\pi$ is $C$-$\sigma$-blocked''. Applying step~2 pointwise inside this universal -- replacing the body ``$\pi$ is $C$-$\sigma$-blocked'' with its simplified disjunction on each walk $\pi$ separately -- yields the equivalent simplified surface: $A \isPerp_G B \given C$ iff every walk from a node in $A$ to a node in $J \cup B$ contains either a collider outside $\Anc^G(C)$ or a (any) non-collider in $C$. This is the LN's promised simplification of $i\sigma$-separation in the acyclic case.

\smallskip
\noindent\textit{Discussion paragraph (informal acknowledgment).}\;
The LN's continuation names the simplified surface obtained in step~3 \emph{$id$-separation}; this naming is formalised separately as a definition in subsection~3.4, with the simplified blocking condition baked directly into the definition rather than gated by an acyclicity hypothesis. The LN's accompanying observations -- that $id$- and $i\sigma$-separation diverge in the non-acyclic setting, that $i\sigma$-separation is the more general concept, and the references listed in the remark above -- motivate that subsection's framing but do not require a separate formal claim in the present row, whose content stops at the acyclic-case simplification just proved.
\end{proof}

\end{document}
